\Article{30}

Methodus alia, summam $W$ generaliter determinandi, petitur ex iis, quae in artt. $22.\ 24$ exposita sunt.\quad Statuamus $\cos\omega+i\sin\omega=\rho$, atque

\[
\begin{aligned}
\rho^{\frac{nn}{aa}}=\alpha,\quad\rho^{\frac{nn}{bb}}=\beta,\quad\rho^{\frac{nn}{cc}}=\gamma\;\text{etc.}
\end{aligned}
\]

ita ut habeatur $r=\rho^k$, $A=\alpha^k$, $B=\beta^k$, $C=\gamma^k$ etc.\quad Tunc erit

\[
\begin{aligned}
1+\rho+\rho^4+\rho^9+\text{etc.}+\rho^{(n-1)^2}
\end{aligned}
\]

productum e factoribus

\[
\begin{aligned}
&1+\alpha+\alpha^4+\alpha^9+\text{etc.}+\alpha^{(a-1)^2}\\
&1+\beta+\beta^4+\beta^9+\text{etc.}+\beta^{(b-1)^2}\\
&1+\gamma+\gamma^4+\gamma^9+\text{etc.}+\gamma^{(c-1)^2},\;\text{etc.}
\end{aligned}
\]

\SourceLeaf{50}

adeoque $W$ productum e factoribus

\[
\begin{aligned}
w&=1+\rho+\rho^4+\rho^9+\text{etc.}+\rho^{(n-1)^2}\\[2pt]
\mathfrak A&=\frac{1+A+A^4+A^9+\text{etc.}+A^{(a-1)^2}}{1+\alpha+\alpha^4+\alpha^9+\text{etc.}+\alpha^{(a-1)^2}}\\[2pt]
\mathfrak B&=\frac{1+B+B^4+B^9+\text{etc.}+B^{(b-1)^2}}{1+\beta+\beta^4+\beta^9+\text{etc.}+\beta^{(b-1)^2}}\\[2pt]
\mathfrak C&=\frac{1+C+C^4+C^9+\text{etc.}+C^{(c-1)^2}}{1+\gamma+\gamma^4+\gamma^9+\text{etc.}+\gamma^{(c-1)^2}}
\end{aligned}
\]

Iam factor primus $w$ determinatus est per disquisitiones supra traditas (art. $19$); factores reliqui vero $\mathfrak A,\mathfrak B,\mathfrak C$ etc. prodeunt per formulas artt. $22.\ 24$, quas ut omnia iuncta habeantur, hic denuo colligimus\begingroup\renewcommand{\thefootnote}{*)}\footnote{Manifesto, quae illic erant $k$ et $h$, hic erunt $\frac{n}{a}$ et $k$ respectu factoris secundi, $\frac{n}{b}$ et $k$ respectu factoris tertii etc.}\endgroup.\quad Duodecim casus hic sunt distinguendi, scilicet

I.\quad Si $a$ est numerus primus (impar) $=p$, vel talis numeri potestas cum exponente impari, atque $k$ residuum quadraticum ipsius $p$, erit factor respondens $\mathfrak A=+1$.

II.\quad Si manentibus reliquis $k$ est non-residuum quadraticum ipsius $p$, erit $\mathfrak A=-1$.

III.\quad Si $a$ est quadratum numeri primi imparis, altiorve eius potestas cum exponente pari, erit $\mathfrak A=+1$.

IV.\quad Si $a$ est $=4$, aut altior binarii potestas cum exponente pari, simulque $k$ formae $4\mu+1$, erit $\mathfrak A=+1$.

V.\quad Si, manentibus reliquis ut in IV, $k$ est formae $4\mu+3$, atque $\frac{n}{a}$ formae $4\mu+1$, erit $\mathfrak A=-i$.

VI.\quad Si, manentibus reliquis ut in IV, $k$ est formae $4\mu+3$, atque $\frac{n}{a}$ formae $4\mu+3$, erit $\mathfrak A=+i$.

VII.\quad Si $a$ est $=8$, aut altior binarii potestas cum exponente impari, atque $k$ formae $8\mu+1$, erit $\mathfrak A=+1$.

VIII.\quad Si, manentibus reliquis ut in VII, $k$ est formae $8\mu+5$, erit $\mathfrak A=-1$.

IX.\quad Si, manentibus reliquis ut in VII, $k$ est formae $8\mu+3$, atque $\frac{n}{a}$ formae $4\mu+1$, erit $\mathfrak A=+i$.

\SourceLeaf{51}

X.\quad Si, manentibus reliquis ut in VII, $k$ est formae $8\mu+3$, atque $\frac{n}{a}$ formae $4\mu+3$, erit $\mathfrak A=-i$.

XI.\quad Si, manentibus reliquis ut in VII, $k$ est formae $8\mu+7$, atque $\frac{n}{a}$ formae $4\mu+1$, erit $\mathfrak A=-i$.

XII.\quad Si, manentibus reliquis ut in VII, $k$ est formae $8\mu+7$, atque $\frac{n}{a}$ formae $4\mu+3$, erit $\mathfrak A=+i$.

Casum eum, ubi $a=2$, praeterimus; hic quidem $\mathfrak A$ foret $=\frac00$ sive indeterminatus, sed tunc semper $W=0$.

Factores reliqui $\mathfrak B,\mathfrak C$ etc. perinde pendent a $b,c$ etc., ut $\mathfrak A$ ab $a$, quatenus in illorum determinationem ingrediuntur.

\Article{31}

Secundum hanc methodum alteram exemplum primum art. $29$ ita se habet:

{\small
\[
\begin{aligned}
&\text{Factor }w\text{ fit }=(1+i)\surd2520\\
&\text{Pro }a=8\text{ factor respondens }\mathfrak A\text{ fit, per casum VIII, }=-1\\
&\text{Factori ipsius }n\text{ secundo }9\text{ respondet factor }+1\;\text{(per casum III.)}\\
&\text{Factori }5\text{ respondet factor }-1\;\text{(per casum II.)}\\
&\text{Factori }7\text{ respondet factor }-1\;\text{(per casum II.)}
\end{aligned}
\]
}

Hinc conflatur productum $W=(-1-i)\surd2520$, ut in art. $29$.

\Article{32}

Quum valor ipsius $W$ per methodos \textit{duas} determinari possit, quarum altera relationibus numerorum $\frac{nk}{a}$, $\frac{nk}{b}$, $\frac{nk}{c}$ etc. ad numeros $a,b,c$ etc. innititur, altera vero a relationibus ipsius $k$ ad numeros $a,b,c$ etc. pendet, inter omnes has relationes nexus quidam conditionalis intercedere debet, ita ut quaevis e reliquis determinabilis esse debeat.\quad Supponamus, omnes numeros $a,b,c$ etc. esse numeros primos impares, atque $k$ accipi $=1$; distribuanturque factores $a,b,c$ etc. in duas classes, quarum altera contineat eos, qui sunt formae $4\mu+1$, et qui denotentur per $p,p',p''$ etc., altera vero constet ex iis, qui sunt formae $4\mu+3$, et qui exprimantur per $q,q',q''$ etc.: multitudinem posteriorum designabimus per $m$. His ita factis, observamus primo, $n$ fieri formae $4\mu+1$, si $m$ fuerit par (quorsum etiam referri debet casus is, ubi factores classis alterius omnino desunt, sive ubi $m=0$), contra $n$ fieri formae $4\mu+3$, si $m$ fuerit impar.\quad Iam determinatio
\SourceLeaf{52}

ipsius $W$ per methodum primam ita perficitur.\quad Pendeant numeri $P,P'\;P''$ etc., $Q,Q',Q''$ etc. ita a relationibus numerorum $\frac{n}{p}$, $\frac{n}{p'}$, $\frac{n}{p''}$ etc., $\frac{n}{q}$, $\frac{n}{q'}$, $\frac{n}{q''}$ etc. ad numeros $p,p',p''$ etc., $q,q',q''$ etc. resp., ut statuatur

\[
\begin{aligned}
P&=+1,\;\text{si }\frac{n}{p}\text{ est residuum quadraticum ipsius }p\\
P&=-1,\;\text{si }\frac{n}{p}\text{ est non-residuum quadraticum ipsius }p
\end{aligned}
\]

et perinde de reliquis.\quad Tunc erit $W$ productum e factoribus $P\surd p$, $P'\surd p'$, $P''\surd p''$ etc. $iQ\surd q$, $iQ'\surd q'$, $iQ''\surd q''$ etc., adeoque

\[
\begin{aligned}
W=PP'P''\ldots QQ'Q''\ldots i^m\surd n
\end{aligned}
\]

Per methodum secundam, aut potius statim per praecepta art. $19$, erit

\[
\begin{aligned}
W&=+\surd n,\;\text{si }n\text{ est formae }4\mu+1,\\&\qquad\text{vel quod eodem redit, si }m\text{ est par}\\
W&=+i\surd n,\;\text{si }n\text{ est formae }4\mu+3,\\&\qquad\text{vel si }m\text{ est impar.}
\end{aligned}
\]

Utrumque casum simul complecti licet per formulam sequentem:

\[
\begin{aligned}
W=i^{mm}\surd n
\end{aligned}
\]

Hinc itaque colligitur

\[
\begin{aligned}
PP'P''\ldots QQ'Q''\ldots=i^{mm-m}
\end{aligned}
\]

Sed $i^{mm-m}$ fit $=1$, quoties $m$ est formae $4\mu$ vel $4\mu+1$, atque $=-1$, quoties $m$ est formae $4\mu+2$ vel $4\mu+3$, unde deducimus sequens elegantissimum

\textsc{Theorema.}\quad \textit{Denotantibus $a,b,c$ etc. numeros primos impares positivos inae}\textit{quales, quorum productum statuitur $=n$, et inter quos $m$ sint formae $4\mu+3$, re}\textit{liqui formae $4\mu+1$: multitudo eorum ex his numeris $a,b,c$ etc., quorum non-resi}\textit{dua resp. sunt $\frac{n}{a}$, $\frac{n}{b}$, $\frac{n}{c}$ etc., par erit, quoties $m$ est formae $4\mu$ vel $4\mu+1$, impar} \textit{vero, quoties $m$ est formae $4\mu+2$ vel $4\mu+3$.}

Ita e. g. statuendo $a=3$, $b=5$, $c=7$, $d=11$, habemus tres numeros formae $4\mu+3$, puta $3$, $7$ et $11$; est autem $5\pdot7\pdot11\,R3$; $3\pdot7\pdot11\,R5$; $3\pdot5\pdot11\,R7$; $3\pdot5\pdot7\,N11$, sive unicus $\frac{n}{d}$ est non-residuum ipsius $d$.

\Article{33}

Celeberrimum \textit{theorema fundamentale} circa residua quadratica nihil aliud est, nisi casus specialis theorematis modo evoluti.\quad Limitando scilicet multitudinem
\SourceLeaf{53}

numerorum $a,b,c$ etc. ad \textit{duos}, patet, si unus tantum ex ipsis, vel neuter, sit formae $4\mu+3$, fieri debere vel simul $aRb$, $bRa$, vel simul $aNb$, $bNa$; contra si uterque est formae $4\mu+3$, unus ex ipsis alterius non-residuum esse debebit, atque hic illius residuum.\quad En itaque demonstrationem \textit{quartam} huius gravissimi theorematis, cuius demonstrationem primam et secundam in Disquisitionibus Arithmeticis, tertiam nuper in commentatione peculiari tradidimus (\textit{Commentt. T. XVI}): duas alias principiis rursus omnino diversis innitentes in posterum exponemus.\quad Summopere sane est mirandum, quod hocce venustissimum theorema, quod primo omnes conatus tam pertinaciter eluserat, tot postea viis toto coelo inter se distantibus adiri potuerit.

\Article{34}

Etiam theoremata reliqua, quae quasi supplementum ad theorema fundamentale efficiunt, scilicet per quae dignoscuntur numeri primi, quorum residua vel non-residua sunt $-1,+2$ et $-2$, ex iisdem principiis derivari possunt. Incipiemus a residuo $+2$.

Statuendo $n=8a$, ita ut $a$ sit numerus primus, atque $k=1$, per methodum art. $28$. $W$ erit productum e duobus factoribus, quorum alter erit $+\surd a$, vel $+i\surd a$, si $8$, vel quod idem est $2$, est residuum quadraticum ipsius $a$; contra $-\surd a$ vel $-i\surd a$, si $2$ est non-residuum ipsius $a$.\quad Factor secundus autem est

\[
\begin{aligned}
&(1+i)\surd8,\quad\text{si }a\text{ est formae }8\mu+1\\
&(-1+i)\surd8,\quad\text{si }a\text{ est formae }8\mu+3\\
&(-1-i)\surd8,\quad\text{si }a\text{ est formae }8\mu+5\\
&(1-i)\surd8,\quad\text{si }a\text{ est formae }8\mu+7
\end{aligned}
\]

Sed per art. $18$ semper erit $W=(1+i)\surd n$; dividendo hunc valorem per quatuor valores factoris secundi, patet, factorem primum fieri debere

\[
\begin{aligned}
&+\surd a,\quad\text{si }a\text{ est formae }8\mu+1\\
&-i\surd a,\quad\text{si }a\text{ est formae }8\mu+3\\
&-\surd a,\quad\text{si }a\text{ est formae }8\mu+5\\
&+i\surd a,\quad\text{si }a\text{ est formae }8\mu+7
\end{aligned}
\]

Hinc sponte sequitur, in casu primo et quarto $2$ esse debere residuum ipsius $a$, in casu secundo et tertio autem non-residuum.

\SourceLeaf{54}

\Article{35}

Numeri primi, quorum residuum vel non-residuum est $-1$, facile dignoscuntur adiumento theorematis sequentis, quod etiam per se ipsum satis memorabile est.

\textsc{Theorema.}\quad \textit{Productum e duobus factoribus}

\[
\begin{aligned}
W'&=1+r^{-1}+r^{-4}+\text{etc.}+r^{-(n-1)^2}\\
W&=1+r+r^4+\text{etc.}+r^{(n-1)^2}
\end{aligned}
\]

\textit{est $=n$, si $n$ est impar; vel $=0$, si $n$ est impariter par; vel $=2n$, si $n$ est} \textit{pariter par.}

\textit{Demonstr.}\quad Quum manifesto fiat

\[
\begin{aligned}
W&=r+r^4+r^9+\text{etc.}+r^{nn}\\
&=r^4+r^9+\text{etc.}+r^{(n+1)^2}\\
&=r^9+\text{etc.}+r^{(n+2)^2}\quad\text{etc.}
\end{aligned}
\]

productum $WW'$ ita quoque exhiberi poterit

\[
\begin{aligned}
&1+r+r^4+r^9+\text{etc.}+r^{(n-1)^2}\\
+{}&r^{-1}(r+r^4+r^9+r^{16}+\text{etc.}+r^{nn})\\
+{}&r^{-4}(r^4+r^9+r^{16}+r^{25}+\text{etc.}+r^{(n+1)^2})\\
+{}&r^{-9}(r^9+r^{16}+r^{25}+r^{36}+\text{etc.}+r^{(n+2)^2})\\
&\text{etc.}\\
+{}&r^{-(n-1)^2}(r^{(n-1)^2}+r^{nn}+r^{(n+1)^2}\\&\qquad+r^{(n+2)^2}+\text{etc.}+r^{(2n-2)^2})
\end{aligned}
\]

quod aggregatum verticaliter summatum producit

\[
\begin{aligned}
&n\\
+{}&r(1+rr+r^4+r^6+\text{etc.}+r^{2n-2})\\
+{}&r^4(1+r^4+r^8+r^{12}+\text{etc.}+r^{4n-4})\\
+{}&r^9(1+r^6+r^{12}+r^{18}+\text{etc.}+r^{6n-6})\\
+{}&\text{etc.}\\
+{}&r^{(n-1)^2}(1+r^{2n-2}+r^{4n-4}\\&\qquad+r^{6n-6}+\text{etc.}+r^{2(n-1)^2})
\end{aligned}
\]

Iam si $n$ impar est, singulae partes huius aggregati, praeter primam $n$, erunt $=0$; secunda enim manifesto fit $\frac{r(1-r^{2n})}{1-rr}$, tertia $\frac{r^4(1-r^{4n})}{1-r^4}$ etc. Quoties vero $n$ par est, excipere insuper oportebit partem

\SourceLeaf{55}

\[
\begin{aligned}
r^{\frac14nn}(1+r^n+r^{2n}+r^{3n}+\text{etc.}+r^{nn-n})
\end{aligned}
\]

quae fit $=nr^{\frac14nn}$. In casu priori itaque fit $WW'=n$, in posteriori autem $=n+nr^{\frac14nn}$; sed $r^{\frac14nn}$ fit $=+1$, si $n$ est pariter par, tunc itaque prodit $WW'=2n$; contra fit $r^{\frac14nn}=-1$, si $n$ est impariter par, ubi itaque evadit $WW'=0$. Q. E. D.

\Article{36}

Iam per art. 22 constat, si $n$ sit numerus primus impar, $\frac{W'}{W}$ fieri $=+1$ vel $=-1$, prout $-1$ fuerit residuum vel non-residuum ipsius $n$. Hinc in casu priori esse debebit $W^2=+n$, in posteriori $W^2=-n$; quamobrem per art. 13 concludimus, casum priorem tunc tantum locum habere posse, quando $n$ sit formae $4\mu+1$, casumque posteriorem, quando $n$ sit formae $4\mu+3$.

Denique e combinatione conditionum pro residuis $+2$ et $-1$ inventarum sponte sequitur, $-2$ esse residuum cuiusvis numeri primi formae $8\mu+1$ vel $8\mu+3$, atque non-residuum cuiusvis numeri primi formae $8\mu+5$ vel $8\mu+7$.

\BookRule{42mm}
